%------------------------------------------------------------------------------
% Template file for the submission of papers to IUCr journals in LaTeX2e
% using the iucr document class
% Copyright 1999-2013 International Union of Crystallography
% Version 1.6 (28 March 2013)
%------------------------------------------------------------------------------

\documentclass[preprint]{iucr}              % DO NOT DELETE THIS LINE

     %-------------------------------------------------------------------------
     % Information about journal to which submitted
     %-------------------------------------------------------------------------
     \journalcode{J}              % Indicate the journal to which submitted
                                  %   A - Acta Crystallographica Section A
                                  %   B - Acta Crystallographica Section B
                                  %   C - Acta Crystallographica Section C
                                  %   D - Acta Crystallographica Section D
                                  %   E - Acta Crystallographica Section E
                                  %   F - Acta Crystallographica Section F
                                  %   J - Journal of Applied Crystallography
                                  %   M - IUCrJ
                                  %   S - Journal of Synchrotron Radiation
\usepackage{graphicx}
\usepackage[T1]{fontenc}
\usepackage[utf8]{inputenc}
\usepackage{amsmath}

\begin{document}                  % DO NOT DELETE THIS LINE

     %-------------------------------------------------------------------------
     % The introductory (header) part of the paper
     %-------------------------------------------------------------------------

     % The title of the paper. Use \shorttitle to indicate an abbreviated title
     % for use in running heads (you will need to uncomment it).

\title{Coping with the parallax effect in 2D X-ray detectors}
%\shorttitle{Short Title}

     % Authors' names and addresses. Use \cauthor for the main (contact) author.
     % Use \author for all other authors. Use \aff for authors' affiliations.
     % Use lower-case letters in square brackets to link authors to their
     % affiliations; if there is only one affiliation address, remove the [a].

\cauthor[a]{Jérôme}{Kieffer}{jerome.kieffer@esrf.fr}{}
\author[a]{Edgar}{Gutierrez-Fernandez}
\author[a]{Thomas}{Vincent} 
\author[a]{Loïc}{Huder} 
\author[b]{Gudrun}{Lotze}
\author[c]{Nils}{Blanc}
\author[c]{victor}{Poline}

\aff[a]{European Synchrotron Radiation Facility;71, avenue des Martyrs;CS 40220;38043 Grenoble Cedex 9 \country{France}}
\aff[b]{TBD}
\aff[c]{Institut Néel CNRS; 25 avenue des Martyrs; BP 166; 38042 Grenoble Cedex 9 \country{France}}
     % Use \shortauthor to indicate an abbreviated author list for use in
     % running heads (you will need to uncomment it).

%\shortauthor{Soape, Author and Doe}

     % Use \vita if required to give biographical details (for authors of
     % invited review papers only). Uncomment it.

%\vita{Author's biography}

     % Keywords (required for Journal of Synchrotron Radiation only)
     % Use the \keyword macro for each word or phrase, e.g. 
     % \keyword{X-ray diffraction}\keyword{muscle}

\keyword{Thick sensor}
\keyword{Diffraction setup calibration}
\keyword{Detector physics}

     % PDB and NDB reference codes for structures referenced in the article and
     % deposited with the Protein Data Bank and Nucleic Acids Database (Acta
     % Crystallographica Section D). Repeat for each separate structure e.g
     % \PDBref[dethiobiotin synthetase]{1byi} \NDBref[d(G$_4$CGC$_4$)]{ad0002}

%\PDBref[optional name]{refcode}
%\NDBref[optional name]{refcode}

\maketitle                        % DO NOT DELETE THIS LINE

\begin{synopsis}
Numerical simulation and correction of the 2θ values due to parallax effect, i.e. volumetric absorption of photons in the thickness of the sensor layer.
\end{synopsis}

\begin{abstract}

The parallax effect in X-ray area detectors arises when X-rays from a point source (e.g. a scattering sample) traverse thick detection layers, causing a shift in the apparent position of detected events. This can degrade the angular resolution and introduce artifacts, especially in pair distribution function (PDF) analysis and powder diffraction measurements.
In thick detectors, X-rays at high scattering angles may pass through multiple pixels, leading to discrepancies between observed and actual scattering angles (2θ)

\end{abstract}


     %-------------------------------------------------------------------------
     % The main body of the paper
     %-------------------------------------------------------------------------
     % Now enter the text of the document in multiple \section's, \subsection's
     % and \subsubsection's as required.

\section{Introduction}

Experimental setups are usually calibrated on synchrotron sources using the powder sample $\rm LaB_6, CeO_2$ diffraction pattern that is properly characterized (by NIST).
Those samples show nice concentric Debye-Scherrer rings in their diffraction pattern that are analyzed with tools such as pyFAI-calib2 \cite{pyfai_2020}, Dawn \cite{fil} or 

\section{State of the art:}
\section{Introduction}


\section{Conclusion}

     %-------------------------------------------------------------------------
     % The back matter of the paper - acknowledgements and references
     %-------------------------------------------------------------------------

     % Acknowledgements come after the appendices

\ack{Acknowledgements:}

D2AM / BM02 

Vincent Favre Nicolin

We would like to thank also Jonathan P. Wright and Carlotta Giacobbe for offering us the ability to test those algorithms on small molecule data and validate the concept of sparsification on the ID11 beamline at the ESRF.

\bibliographystyle{iucr}
\bibliography{biblio}

\end{document}                    % DO NOT DELETE THIS LINE
%%%%%%%%%%%%%%%%%%%%%%%%%%%%%%%%%%%%%%%%%%%%%%%%%%%%%%%%%%%%%%%%%%%%%%%%%%%%%%